If the digit product is positive and $n$ does not end in nine, changing its final digit to the next digit changes the product. If it ends in nine, $n+1$ has a zero digit and its product is zero. Thus equality requires both products to be zero.
This happens exactly when the decimal prefix before the last digit of $n$ contains zero. If no carry occurs, that zero remains. If a carry occurs through trailing nines, the new trailing zeros ensure the next product is zero. Conversely, a zero only in the units position disappears on adding one, so it is insufficient.
For three-digit numbers, the two-digit prefix has a zero in $9$ cases, giving $90$ numbers. For four-digit numbers from $1000$ through $1999$, the prefix $100$ through $199$ contains zero in $100-81=19$ cases, giving $190$. All twelve numbers from $2000$ through $2011$ qualify. No one- or two-digit number does. The answer is $\boxed{90+190+12=292}$.
